%A LaTeX model for SYSU General Phys. Lab. by Probfia Gao.
%用XeLaTeX编译
\documentclass[11pt,a4paper]{ctexart}

%在下面补全实验名，例如 实验BB3 光电效应实验。
\newcommand{\ExpeName}{实验D5 CCD原理与应用}

\usepackage{fancyhdr}
\usepackage{amsmath}
\usepackage{graphicx}
\usepackage[hmargin=1.25in,vmargin=1in]{geometry}
\usepackage{pdfpages}
\usepackage[colorlinks,
            linkcolor=red,
		 urlcolor=black]{hyperref}
\usepackage{cleveref}
\usepackage{float}


\crefname{equation}{}{}
\crefname{figure}{图}{图}
\crefname{footnote}{注释}{注释}
\crefname{table}{表}{表}

%\cpic{<尺寸>}{<文件名>}}用于生成居中的图片。
\newcommand{\cpic}[2]{
\begin{center}
\includegraphics[scale=#1]{#2}
\end{center}
}

%\cpicn{<尺寸>}{<文件名>}{<注释>}用于生成居中且带有注释的图片，其label为图片名。
\newcommand{\cpicn}[3]
{
\begin{figure}[H]
\cpic{#1}{#2}
\caption{#3\label{#2}}
\end{figure}
}

\newcommand{\beq}{\begin{equation}}
\newcommand{\eeq}{\end{equation}}
\newcommand{\bea}{\begin{equation}\begin{aligned}}
\newcommand{\eea}{\end{aligned}\end{equation}}

%输入单位和数学常数
%下面所有命令需在公式环境下使用
\newcommand{\e}{\mathrm{e}}   %自然常数e = \e
\newcommand{\im}{\mathrm{i}}   %虚数单位i = \im
\newcommand{\meter}{\mathrm{m}}      %单位/前缀 = \单位/前缀英文名
\newcommand{\newton}{\mathrm{N}}  
\newcommand{\joule}{\mathrm{J}}
\newcommand{\second}{\mathrm{s}}
\newcommand{\gram}{\mathrm{g}}
\newcommand{\ampere}{\mathrm{A}}
\newcommand{\kilo}{\mathrm{k}}
\newcommand{\milli}{\mathrm{m}}
\newcommand{\kelvin}{\mathrm{K}}
\newcommand{\mole}{\mathrm{mol}}
\newcommand{\volt}{\mathrm{V}}
\newcommand{\nano}{\mathrm{n}}
\newcommand{\degreeC}{^\circ \mathrm{C}}  %摄氏度符号 = \degreeC

\newcommand{\unit}[1]{\rm \ #1}
\newcommand{\emptyline}{\par \ \\ }

\pagestyle{fancy}

\fancyhead[L]{\footnotesize{中山大学物理与天文学院近代物理实验}}
\fancyhead[R]{\footnotesize{\ExpeName}}
\fancyfoot[C]{\thepage}

\begin{document}
%第一页
\cpic{0.255}{e1}%学生信息和计分表格
\begin{table}[H]
\centering
\begin{tabular}{|p{32mm}|p{32mm}|p{32mm}|p{32mm}|}
\hline
年级、专业： & 17级 物理学 & 组号： & 实验班6 \\ \hline
姓名： & 高寒 & 学号： & 17353019 \\ \hline
日期： & \today & 教师签名： &  \\ \hline
\end{tabular}
\end{table}
\begin{center}
\LARGE\textbf{{\ExpeName}}
\end{center}
\large{【实验报告注意事项】}
\begin{enumerate}
 \item 实验报告由三部分组成：
 \begin{enumerate}
  \item[1)]预习报告：（提前一周）认真研读\textbf{\uline{实验讲义}}，弄清实验原理；实验所需的仪器设备、用具及其使用（强烈建议到实验室预习），完成讲义中的预习思考题；了解实验需要测量的物理量，并根据要求提前准备实验记录表格（由学生自己在实验前设计好，可以打印）。预习成绩低于50\%者不能做实验{\color{red} (实验D2和D3需要提前一周的周四完成预习报告交任课老师批改，批改通过后，才允许做实验)}。

  \item[2)]实验记录：认真、客观记录实验条件、实验过程中的现象以及数据。实验记录请用珠笔或者钢笔书写并签名（{\color{red}用铅笔记录的被认为无效}）。{\color{red}保持原始记录，包括写错删除部分，如因误记需要修改记录，必须按规范修改。}（不得输入电脑打印，但可扫描手记后打印扫描件）；离开前请实验教师检查记录并签名。
  \item[3)]分析讨论：处理实验原始数据（学习仪器使用类型的实验除外），对数据的可靠性和合理性进行分析；按规范呈现数据和结果（图、表），包括数据、图表按顺序编号及其引用；分析物理现象（含回答实验思考题，写出问题思考过程，必要时按规范引用数据）；最后得出结论。
 \end{enumerate}
 \textbf{实验报告}就是预习报告、实验记录、和数据处理与分析合起来，加上本页封面。
 \item 每次完成实验后的一周内交\textbf{实验报告}。
 \item 除实验记录外，实验报告其他部分建议双面打印。
\end{enumerate}
\ 
\\
\ 

\begin{flushright}                                                           %模板作者
\tiny{
A \LaTeX \ model for Modern Phys. Lab., SPA, SYSU by \em{\href{https://www.weibo.com/3532532974/profile?rightmod=1&wvr=6&mod=personinfo&is_all=1}{Probfia} Gao.}
}
\end{flushright}
\includepdf[pages=2-14]{preview.pdf}
\iffalse
\newpage%预习报告
\begin{center}
\LARGE{\textbf{\ExpeName}}
\end{center}
\textbf{【实验目的】}
\begin{enumerate}
 \item[1.] (目的1)
 \item[2.] (目的2)
\end{enumerate}
\textbf{【仪器用具】}
%将讲义中的表格截图保存为t1在该文件夹下后删去下一行之前的%符号，合理调整scale参数。
%cpic{0.3}{t1}
%或者自己去 https://www.tablesgenerator.com/ 做一个表。
\textbf{【原理概述】}\par
该实验\\
\textbf{【实验前思考题】}
\begin{enumerate}
 \item[1.] (问题1)
 \item[2.] (问题2)
\end{enumerate}

\newpage%实验记录
%\cpic{0.255}{e2}%学生信息表格
\begin{table}[H]
\centering
\begin{tabular}{|p{32mm}|p{32mm}|p{32mm}|p{32mm}|}
\hline
年级、专业： & 17级 物理学 & 组号： & 实验班6 \\ \hline
姓名： & 高寒 & 学号： & 17353019 \\ \hline
日期： &  & 实验地点： & 珠海教学楼 A102 \\ \hline
学生签名： &  
\includegraphics[scale=0.09]{sign}& 评分： &  \\ \hline
日期： &  & 教师签名： &  \\ \hline
\end{tabular}
\end{table}

\begin{center}
\LARGE{\textbf{\ExpeName}}
\end{center}
\textbf{【实验内容、步骤、结果】}
\par
(表格)
\newline
\textbf{【实验过程中遇到问题记录】}\par
\cpic{0.81}{none}
\fi
%生成最终报告时将上面内容全部删除或注释（用\iffalse \fi），将扫描得到的实验报告保存为Record.pdf在LaTeX model for GPL下，将下行命令的注释号删去。注意根据实际页数调整pages参数。
\includepdf[pages=1-9]{Record}


\newpage%分析与讨论
%\cpic{0.255}{e3}%学生信息表格
\begin{table}[H]
\centering
\begin{tabular}{|p{32mm}|p{32mm}|p{32mm}|p{32mm}|}
\hline
年级、专业： & 17级 物理学 & 组号： & 实验班6 \\ \hline
姓名： & 高寒 & 学号： & 17353019 \\ \hline
日期： & \today &  &  \\ \hline
评分： &  & 教师签名： &  \\ \hline
\end{tabular}
\end{table}
\begin{center}
\LARGE\textbf{{\ExpeName}}
\end{center}
\textbf{【分析与讨论】}\\
1.\textbf{线阵CCD驱动及特性分析}\par
按要求接线，连接示波器\cite{dtem}发现，SH和RS的输出如\cref{RS}
\cpicn{0.08}{RS}{复位信号RS的波形}
周期分别为
\bea
T_{SH} &= 7.0 \unit{ms} \\
T_{RH} &= 1.4 \unit{\mu s}
\eea
可以发现两者之间的关系为
\beq
T_{SH} = 5000 T_{RH}
\eeq
与工作时序图的示意相同。\par
\cpicn{0.08}{revp}{$\Phi_1$和$\Phi_2$的反相输出}
如\cref{revp}发现$\Phi_1$和$\Phi_2$的周期相同，为
\beq
T_{\Phi_1} = T_{\Phi_2} = 2.780 \unit{\mu s}
\eeq
且相位相差$180^\circ$，两者反相\cite{ced}交替变化，使得模拟移位寄存器在$\Phi_1$与$\Phi_2$脉冲的作用下驱使$\Phi_1$ 电极下势阱中的信号电荷向左移动，并经输出电路由 OS 电极输出。\cite{lec,att1,ccdwiki}
\par
观察CCD端在光照和遮挡下的输出如\cref{light,cover}，发现完全遮挡下，CCD不再输出脉冲信号，与预期相符。
\cpicn{0.7}{light}{CCD输出端在光照下的输出信号}
\cpicn{0.08}{cover}{CCD输出端在镜头被遮挡时的输出信号}
\par
不同阈值电压下的CCD二值化输出如\cref{hvolt,lvolt}。
\cpicn{0.08}{hvolt}{高阈值电压下的CCD二值化输出}
\cpicn{0.08}{lvolt}{低阈值电压下的CCD二值化输出}
可以看到，低阈值电压下波包非常明显，体现出对外界的更灵敏响应。
\par
连接电脑，遮挡不同部位，观察到波形变化如\cref{open,close,l43,r12}。
\cpicn{0.8}{open}{不遮挡时CCD输出波形}
\cpicn{0.8}{close}{全遮挡时CCD输出波形}
\cpicn{0.8}{l43}{遮挡左边3/4时CCD的输出波形}
%\cpicn{0.8}{r12}{遮挡右半部分的CCD输出波形}
观察可以发现，显示屏上的波形与线阵CCD上某点接受到的光强相关。
\par
加入遮挡物，修改CCD频率，输出波形变化如\cref{f300,f358,f444,f600,f800}。
\cpicn{0.8}{f300}{CCD相机频率为$300\unit {kHz}$时的输出波形}
\cpicn{0.8}{f358}{CCD相机频率为$358\unit {kHz}$时的输出波形}
\cpicn{0.8}{f444}{CCD相机频率为$444\unit {kHz}$时的输出波形}
\cpicn{0.8}{f600}{CCD相机频率为$600 \unit {kHz}$时的输出波形}
\cpicn{0.8}{f800}{CCD相机频率为$800\unit {kHz}$时的输出波形}
可以在波形上明显看出被测物的遮挡作用，且频率越大，遮挡越明显。
\par
调整积分时间，波形变化如\cref{int4999,int7999,int12499,int16199,int19999}。
\cpicn{0.8}{int4999}{积分时间为$4999 \unit{\mu s}$的输出波形}
\cpicn{0.8}{int7999}{积分时间为$7999 \unit{\mu s}$的输出波形}
\cpicn{0.8}{int12499}{积分时间为$12499 \unit{\mu s}$的输出波形}
\cpicn{0.8}{int16199}{积分时间为$16199 \unit{\mu s}$的输出波形}
\cpicn{0.8}{int19999}{积分时间为$19999 \unit{\mu s}$的输出波形}
可以看到，在高积分时间下，遮挡物的作用越加不明显，波形趋于饱和。
\emptyline
\\
2.\textbf{利用线阵CCD测量工作直径和角位置}\par
根据之前的实验发现，驱动频率为$800 \unit{kHz}$，脉冲积分时间为$9999 \unit{\mu s}$时，被测物的遮挡作用较为显著。利用$\phi 12$的标准器件得到实际物理尺寸与像素值的换算关系为
\beq
12\unit{mm} \equiv 842.654478 \unit{pix}
\eeq
测量物体直径如\cref{12}
\cpicn{0.8}{12}{测量$\phi 12$物体的直径}
四次重复实验的结果为
\bea
d_1 &= 11.931166\unit{mm} \\
d_2 &= 11.954111\unit{mm} \\
d_3 &= 12.015432 \unit{mm} \\
d_4 &= 11.975548 \unit{mm}
\eea
平均值
\beq
\bar{d} =\frac{11.931166+11.954111+12.015432+11.975548}{4}\unit{mm}  = 11.969\unit{mm}
\eeq
平均值的标准偏差
\beq
\sigma_{\bar{d}} = 0.036\unit{mm}
\eeq
标准值$12\unit{mm}$落在$1\sigma$范围内，结果较为精确。
\par
类似地测量$\phi10$圆柱体得到
\bea
d_1 &= 10.035462\unit{mm} \\
d_2 &= 9.932675 \unit{mm}\\
d_3 &= 9.989101\unit{mm} \\
d_4 &= 10.065479\unit{mm}
\eea
平均值和标准差为\cite{error}
\beq
d = 10.006 \pm 0.06 \unit{mm}
\eeq
标准值亦落在$1\sigma$误差范围内。
\par
进行角位置测量如\cref{angle}。
\cpicn{0.8}{angle}{角位置测量示意}
测量移动平面镜前后峰的位置变化为$s \unit{pix}$，而反射镜到CCD的距离为$d = 17.85\unit{cm}$，根据换算关系$1\unit{pix} \equiv 14\unit{\mu m}$可知，$d = 12750\unit{pix}$。于是转动的弧度值
\beq
\alpha = \frac{1}{2} \frac{s\unit{pix}}{d} = \frac{s}{25500} \unit{rad}
\eeq
四次转动对应的角度为
\bea
\alpha_1 &= \frac{1430.13-1133.67}{25500} \unit{rad} = 0.01163 \unit{rad}\\
\alpha_2 &= \frac{1165.89-898.44}{25500} \unit{rad} = 0.01049\unit{rad}\\
\alpha_3 &= \frac{1467.19-853.32}{25500} \unit{rad} = 0.02407\unit{rad}\\
\alpha_4 &= \frac{1533.25-1173.95}{25500} \unit{rad} = 0.01409\unit{rad}
\eea
可以发现，该测量方法能够精确地测量小角度角位移。误差可能主要来自峰的弥散。
\emptyline
\\
3.\textbf{面阵CCD图像采集、运算和颜色处理}\par
打开面阵CCD测量软件\cite{att2}，图像清晰后停止捕获，如\cref{image}。
\cpicn{0.6}{image}{捕获的原始图像}
灰度直方图如\cref{grey_hist}
\cpicn{0.8}{grey_hist}{原始图像的灰度直方图}
对图像进行灰度线性变换
\beq
y = 3x
\eeq
如\cref{3x}
\cpicn{0.8}{3x}{进行灰度线性变换$y=3x$下的图像}
可以看出，灰度增大下，图像的亮度变高。\par
进行灰度阈值变换如\cref{Snipaste_2019-09-16_17-51-41,Snipaste_2019-09-16_17-51-56}。
\cpicn{0.8}{Snipaste_2019-09-16_17-51-41}{灰度阈值变换}
\cpicn{0.8}{Snipaste_2019-09-16_17-51-56}{灰度阈值变换}
可以发现，图中的几个阈值变换对图像没有明显改变。
\par
进行灰度拉伸变换后的结果如\cref{Snipaste_2019-09-16_17-51-29}，可以发现图像颜色出现明显畸变、变黑，这是由于灰度过于分散。
\cpicn{0.8}{Snipaste_2019-09-16_17-51-29}{灰度拉伸变换}
而对原图进行灰度均衡变换得到的结果如\cref{Snipaste_2019-09-16_17-53-30}，此时的灰度直方图如\cref{Snipaste_2019-09-16_17-53-37}。
\cpicn{0.8}{Snipaste_2019-09-16_17-53-30}{灰度均衡变换后的图像}
\cpicn{0.8}{Snipaste_2019-09-16_17-53-37}{灰度均衡变换后的灰度直方图}
与\cref{grey_hist}对比可以看到，灰度均衡变换使得灰度分布更加均匀，从而使得图像在视觉上显得更加明亮。
\emptyline
\\
4.\textbf{利用面阵CCD测量工件二维尺寸}\par
利用面阵CCD测量工件圆孔内径如\cref{Snipaste_2019-09-16_17-58-33,Snipaste_2019-09-16_18-01-13,Snipaste_2019-09-16_18-01-24}
\cpicn{0.8}{Snipaste_2019-09-16_17-58-33}{测量工件圆孔内径}
\cpicn{0.8}{Snipaste_2019-09-16_18-01-13}{测量工件圆孔内径}
\cpicn{0.8}{Snipaste_2019-09-16_18-01-24}{测量工件圆孔内径}
结果为
\bea
d_1 &= 2.654\unit{mm} \\
d_2 &= 2.699 \unit{mm} \\
d_3 &= 2.716 \unit{mm}
\eea
均值
\beq
\bar d = \frac{2.654+2.699+2.716}{3}  = 2.690 \unit{(mm)}
\eeq
标准差
\beq
\sigma_d = \frac{1}{\sqrt{3}} \sqrt{\frac{1}{3-1}[(2.654-2.69)^2 +(2.699-2.69)^2 +(2.716-2.69)^2 ]} = 0.018\unit{(mm)}
\eeq
游标卡尺的测量结果为
\beq
d_{ref} = 2.70 \unit{mm}
\eeq
相差无几，说明CCD测量比较精准。
\par
测量圆孔孔距如\cref{Snipaste_2019-09-16_17-59-23,Snipaste_2019-09-16_18-02-21,Snipaste_2019-09-16_18-02-34}
\cpicn{0.8}{Snipaste_2019-09-16_17-59-23}{圆孔孔距测量}
\cpicn{0.8}{Snipaste_2019-09-16_18-02-21}{圆孔孔距测量}
\cpicn{0.8}{Snipaste_2019-09-16_18-02-34}{圆孔孔距测量}
结果为
\bea
d_1 &= 17.929\unit{mm} \\
d_2 &= 17.968 \unit{mm} \\
d_3 &= 18.126 \unit{mm}
\eea
均值
\beq
\bar d = \frac{17.929+17.968+18.126}{3}  = 18.008 \unit{(mm)}
\eeq
标准差
\beq
\sigma_d = \frac{1}{\sqrt{3}} \sqrt{\frac{1}{3-1}[(17.929-18.008)^2 +(17.968-18.008)^2 +(18.126-18.008)^2 ]} = 0.060\unit{(mm)}
\eeq
游标卡尺的测量结果为
\beq
d_{ref} = 18.00 \unit{mm}
\eeq
相差无几，说明CCD测量比较精准。
\par
测量圆孔孔距如\cref{Snipaste_2019-09-16_18-00-41,Snipaste_2019-09-16_18-01-01,Snipaste_2019-09-16_18-01-13}
\cpicn{0.8}{Snipaste_2019-09-16_18-00-41}{平行线距测量}
\cpicn{0.8}{Snipaste_2019-09-16_18-01-01}{平行线距测量}
\cpicn{0.8}{Snipaste_2019-09-16_18-01-13}{平行线距测量}
结果为
\bea
d_1 &= 9.402\unit{mm} \\
d_2 &= 9.407 \unit{mm} \\
d_3 &= 9.405 \unit{mm}
\eea
均值
\beq
\bar d = \frac{9.402+9.407+9.405}{3}  = 9.405 \unit{(mm)}
\eeq
标准差
\beq
\sigma_d = \frac{1}{\sqrt{3}} \sqrt{\frac{1}{3-1}[(9.402-9.405)^2 +(9.407-9.405^2 +(9.405-9.405)^2 ]} = 0.002\unit{(mm)}
\eeq
游标卡尺的测量结果为
\beq
d_{ref} = 9.22 \unit{mm}
\eeq
两者间有一定差距，误差可能源自程序对线状边缘的识别不是很精确。
\par
如果被测物较厚，该方法测量的误差可能较大，因为标定尺明显远于被测物表面使得标定存在系统误差。

\bibliographystyle{IEEEtran}
\bibliography{model.bib}

\end{document}